% ---------------------------------------------------------------------------
%  StressID small-scale pipeline -- results report
%  Compile:  pdflatex results_report.tex   (run twice for the ToC)
%  Only standard packages; no external files required except the optional
%  figure results/small/figures.png (guarded by \IfFileExists).
% ---------------------------------------------------------------------------
\documentclass[11pt,a4paper]{article}

\usepackage[margin=2.3cm]{geometry}
\usepackage{booktabs}
\usepackage{graphicx}
\usepackage{amsmath}
\usepackage{xcolor}
\usepackage[colorlinks=true,linkcolor=black,urlcolor=blue]{hyperref}
\usepackage{fancyvrb}

\newcommand{\code}[1]{\texttt{\small #1}}
\definecolor{alertred}{RGB}{170,30,30}
\newcommand{\alert}[1]{\textcolor{alertred}{\textbf{#1}}}

\title{\textbf{Temporal Cross-Modal Stress Detection on StressID}\\
       \large Small-Scale Pipeline Implementation --- Results Report}
\author{Small-subset run: 16 subjects, 6 tasks, 96 recordings}
\date{28 July 2026}

\begin{document}
\maketitle

\begin{abstract}
\noindent
This report documents a runnable implementation of the StressID research pipeline
(Stages 0--6) and the results of its first full training run on a deliberately small
subset of the dataset. The pipeline executes end to end and the evaluation suite is
complete. \alert{The headline finding is negative and methodological:} a
classifier given \emph{only the three modality-presence bits}, with no signal content
whatsoever, achieves $0.781$ accuracy and $0.774$ macro F1 under leakage-free
GroupKFold --- beating every genuine multimodal baseline. Once modality availability is
held constant, both the temporal and static models collapse to the majority-class
reference. The apparent multimodal performance on this dataset is substantially a
shortcut, not stress signal. This must be resolved before the missing-modality
robustness claim (E12) can be made.
\end{abstract}

\tableofcontents
\vspace{1em}

% ===========================================================================
\section{Scope: what was built and run}

\subsection*{Implemented and executed}
\begin{itemize}\itemsep2pt
  \item \textbf{Stage 0} --- subject-level GroupKFold protocol with leakage assertions;
        corrected classical baselines (E1).
  \item \textbf{Stage 1} --- temporal windowing (10\,s windows, 5\,s hop) with caching.
  \item \textbf{Stage 2} --- modality-specific encoders (physio 1D-CNN+BiLSTM;
        audio log-mel CNN; video per-frame CNN + Transformer).
  \item \textbf{Stage 3} --- multi-head cross-modal attention fusion with modality
        dropout (25\,\% per modality) and CLS aggregation.
  \item \textbf{Stage 4} --- Transformer over the window sequence with task-type
        embedding conditioning.
  \item \textbf{Stage 5} --- multi-task heads (binary stress, 3-class affect,
        stress-score regression) with a contrastive subject-invariance loss.
  \item \textbf{Stage 6} --- evaluation suite: E2, E4, E5, E6, E8, E12, E13, plus a
        shortcut audit added during this work.
\end{itemize}

\subsection*{Not implemented (data or code unavailable in this repository)}
\begin{itemize}\itemsep2pt
  \item \textbf{Stage $-1$} --- multi-dataset pretraining. Requires WESAD, SWELL-KW,
        K-EmoCon, MuSE, ForDigitStress, BESST, EmoWork.
  \item \textbf{Stage 5.5 / E0} --- re-implementation of the three competitor papers.
  \item \textbf{E7, E10, E11} --- cross-dataset transfer and data-scaling ablations;
        all require the auxiliary corpora above.
\end{itemize}

\noindent
Consequently the encoders were trained \emph{from scratch} on roughly 57 recordings per
fold. Stage $-1$ exists in the research plan precisely to address this, and its absence
is the most likely cause of the null result reported in Section~\ref{sec:results}.

% ===========================================================================
\section{Dataset structure}

\subsection{Full StressID corpus}

\begin{Verbatim}[fontsize=\small,frame=single,framesep=3mm]
StressID Dataset/
|-- Physiological/          65 subject folders
|   `-- <subj>/<subj>_<task>.txt    CSV, header "ECG,EDA,RR", 500 Hz
|-- Audio/                  54 subject folders
|   `-- <subj>/<subj>_<task>.wav    16 kHz, mono, 16-bit PCM
|-- Videos/                 53 subject folders
|   `-- <subj>/<subj>_<task>.mp4    5 fps, 1280x720
|-- labels.csv              700 rows: subject/task, binary-stress, affect3-class
|-- labels_supplementary.csv
|-- self_assessments.csv    ';'-separated; tasks as rows, subjects as columns
|-- demographics.csv
`-- Technical information.pdf
\end{Verbatim}

\noindent
64 subjects carry labels, over 11 tasks, for 700 labelled recordings. Modality coverage
is incomplete by design: 12 labelled subjects have no audio directory and 12 have no
video directory.

\begin{table}[h]\centering\small
\caption{Task inventory. Audio exists only for the seven speech tasks.}
\begin{tabular}{lccc}
\toprule
Task & Duration (s) & Audio present & $P(\text{stress})$ (full corpus) \\
\midrule
Breathing  & 180 & no  & 0.172 \\
Relax      & 150 & no  & 0.129 \\
Video1     & 174 & no  & 0.286 \\
Video2     & 118 & no  & 0.397 \\
Reading    & 60  & yes & 0.547 \\
Math       & 60  & yes & 0.672 \\
Speaking   & 60  & yes & 0.672 \\
Counting1  & 60  & yes & 0.656 \\
Counting2  & 60  & yes & 0.766 \\
Counting3  & 60  & yes & 0.703 \\
Stroop     & 60  & yes & 0.766 \\
\bottomrule
\end{tabular}
\end{table}

\noindent
\alert{Note the structural confound visible in this table:} audio is present exactly for
the high-stress tasks. Across the whole corpus,
$P(\text{stress}\mid\text{audio exists}) = 0.683$ versus
$P(\text{stress}\mid\text{no audio}) = 0.246$. Section~\ref{sec:shortcut} quantifies the
consequence.

\subsection{Small subset used in this run}

Subject selection is deterministic and seeded (\code{Config.seed = 1337}): subjects with
complete modality coverage are ranked by how balanced their own binary labels are, then
drawn without replacement.

\begin{center}\small
\begin{tabular}{ll}
\toprule
Subjects & 16 (of 64) --- all with physio, audio and video \\
Subject IDs & \code{2hpu 4woj 6g6y 71i5 8g4y 9t6n 9txq cxj0} \\
            & \code{kkf5 m8g5 p9i3 r5s8 v8mh w2t5 x1q3 y9z6} \\
Tasks & Breathing, Relax, Reading, Math, Stroop, Speaking \\
Recordings & 96 (of 700) \\
Binary balance & 0.531 \\
3-class affect & class 0: 36, class 1: 24, class 2: 36 \\
Cache size & 27.8 MB \\
\bottomrule
\end{tabular}
\end{center}

\begin{table}[h]\centering\small
\caption{Per-task composition of the subset. \code{Windows} is after the
$\text{max\_windows}=16$ cap.}
\begin{tabular}{lccccc}
\toprule
Task & $n$ & $P(\text{stress})$ & Audio & Duration (s) & Windows \\
\midrule
Breathing & 16 & 0.125 & 0.00 & 180 & 16 (truncated) \\
Relax     & 16 & 0.125 & 0.00 & 150 & 16 (truncated) \\
Reading   & 16 & 0.625 & 1.00 & 60  & 11 \\
Math      & 16 & 0.750 & 1.00 & 60  & 11 \\
Speaking  & 16 & 0.688 & 1.00 & 60  & 11 \\
Stroop    & 16 & 0.875 & 1.00 & 60  & 11 \\
\bottomrule
\end{tabular}
\end{table}

\noindent
Breathing and Relax exceed the 16-window cap and are truncated to their first
$\approx 85$\,s. This does not affect E5/E6/E12, which are scored only on the four 60\,s
speech tasks.

\subsection{Cached tensor structure}

Each recording becomes one \code{.npz} under \code{research\_way/data\_small/cache/}:

\begin{Verbatim}[fontsize=\small,frame=single,framesep=3mm]
<subject>_<task>.npz
  physio     float32  [16, 640, 3]        16 windows x 640 samples (64 Hz) x ECG,EDA,RR
  audio      float16  [16, 400,  40]      16 windows x 400 frames x 40 log-mel bins
  video      float16  [16,  25, 32, 32]   16 windows x 25 frames x 32x32 grayscale face
  mask       float32  [16, 3]             per-window modality availability
  n_windows  int64                        valid windows before padding
\end{Verbatim}

\noindent
\textbf{Windowing.} The window grid is defined on the physiological signal --- the only
always-present modality --- and window boundaries are converted to seconds so audio and
video are time-aligned to the same grid.

\noindent
\textbf{Normalisation.} Physio is z-scored \emph{per subject} using statistics pooled over
that subject's own recordings, then clipped to $\pm 10$. Per-recording normalisation was
rejected because it erases the between-task EDA and heart-rate level shifts that carry
most of the stress signal. Audio and video standardisation statistics are computed from
\emph{training subjects only} in each fold.

% ===========================================================================
\section{Evaluation protocol}

Subject-level \code{GroupKFold} with $K=4$. No subject appears in more than one split;
a validation set is carved from training subjects only. Assertions verify train/val/test
disjointness and that every subject is tested exactly once.

\begin{center}\small
\begin{tabular}{lllc}
\toprule
Fold & Test subjects & Test $n$ & Test $P(\text{stress})$ \\
\midrule
0 & \code{71i5 cxj0 r5s8 y9z6} & 24 & 0.50 \\
1 & \code{6g6y 9txq p9i3 x1q3} & 24 & 0.58 \\
2 & \code{4woj 9t6n m8g5 w2t5} & 24 & 0.50 \\
3 & \code{2hpu 8g4y kkf5 v8mh} & 24 & 0.54 \\
\bottomrule
\end{tabular}
\end{center}

\noindent
Training used 3 seeds $\times$ 4 folds $\times$ 2 variants (temporal, static) $=$ 24 runs,
40 epochs each, $\approx$ 20 minutes total on a Quadro P1000. All confidence intervals
below are bootstrap 95\,\% CIs over the 12 (seed $\times$ fold) runs; temporal-vs-static
comparisons use a paired bootstrap.

% ===========================================================================
\section{Results}\label{sec:results}

\subsection{E1 --- corrected baselines under GroupKFold}

\begin{table}[h]\centering\small
\caption{Classical static-feature baselines, leakage-free GroupKFold.
\code{availability\_only} uses \emph{only} the three modality-presence bits and no signal
content at all.}
\begin{tabular}{llcc}
\toprule
Classifier & Features & Accuracy & Macro F1 \\
\midrule
LogReg & physio            & 0.646 & 0.642 \\
LogReg & audio             & 0.708 & 0.704 \\
LogReg & video             & 0.510 & 0.503 \\
LogReg & feature fusion    & 0.708 & 0.705 \\
LogReg & decision fusion   & 0.708 & 0.705 \\
LogReg & \alert{availability only} & \alert{0.781} & \alert{0.774} \\
\midrule
RF & physio            & 0.677 & 0.667 \\
RF & audio             & 0.677 & 0.670 \\
RF & video             & 0.583 & 0.576 \\
RF & feature fusion    & 0.750 & 0.747 \\
RF & decision fusion   & 0.688 & 0.678 \\
RF & \alert{availability only} & \alert{0.781} & \alert{0.774} \\
\bottomrule
\end{tabular}
\end{table}

\begin{table}[h]\centering\small
\caption{Leakage inflation: random (subject-ignoring) KFold minus GroupKFold. Positive
values indicate the leaky protocol overstates performance.}
\begin{tabular}{llcc}
\toprule
Classifier & Features & $\Delta$ Accuracy & $\Delta$ Macro F1 \\
\midrule
LogReg & physio          & $+0.062$ & $+0.048$ \\
LogReg & audio           & $+0.031$ & $+0.005$ \\
LogReg & video           & $+0.021$ & $-0.002$ \\
LogReg & feature fusion  & $-0.021$ & $-0.055$ \\
LogReg & decision fusion & $\phantom{+}0.000$ & $-0.037$ \\
RF & physio          & $+0.010$ & $+0.003$ \\
RF & audio           & $+0.062$ & $+0.041$ \\
RF & video           & $+0.021$ & $-0.004$ \\
RF & feature fusion  & $+0.010$ & $-0.019$ \\
RF & decision fusion & $+0.052$ & $+0.038$ \\
\bottomrule
\end{tabular}
\end{table}

\noindent
At 16 subjects the leakage effect is small and noisy, several entries being negative.
This does \emph{not} refute the leakage narrative --- it reflects that with four test
subjects per fold, fold-to-fold variance dominates. The inflation measurement should be
repeated at full scale before being reported as a contribution.

\subsection{Shortcut audit --- the headline result}\label{sec:shortcut}

The same full-modality predictions are scored two ways. \textbf{``All''} covers every test
recording, where modality availability varies and is label-correlated. \textbf{``Complete''}
covers only recordings where all three modalities exist, where availability is constant
and the shortcut is unavailable.

\begin{table}[h]\centering\small
\caption{Shortcut audit. A model genuinely reading physiology, audio and video would keep
its margin over the majority-class reference in \emph{both} scopes.}
\begin{tabular}{llccccc}
\toprule
Scope & Variant & $n$ & Pos.\ rate & Macro F1 [95\,\% CI] & Majority F1 & Margin \\
\midrule
all      & static   & 96 & 0.531 & 0.621 [0.559, 0.681] & 0.347 & $+0.274$ \\
all      & temporal & 96 & 0.531 & 0.705 [0.650, 0.758] & 0.347 & $+0.358$ \\
complete & static   & 64 & 0.734 & 0.452 [0.388, 0.518] & 0.423 & $+0.028$ \\
complete & temporal & 64 & 0.734 & 0.417 [0.389, 0.450] & 0.423 & \alert{$-0.007$} \\
\bottomrule
\end{tabular}
\end{table}

\noindent
Scored on all recordings this looks like a clean positive result: the temporal model beats
the static model by $0.084$ macro F1, precisely the story E2 is designed to tell. Scored
on the 64 recordings where availability carries no information, \alert{both variants sit
at the majority-class reference} and the temporal model is marginally below it. The
margin was the shortcut.

\subsection{E2 / E4 --- temporal versus static, full modalities}

\begin{center}\small
\begin{tabular}{lccccc}
\toprule
Metric & Temporal & Static & $\Delta$ & 95\,\% CI & $p$ \\
\midrule
Macro F1 & 0.417 & 0.452 & $-0.035$ & $[-0.105, +0.046]$ & 0.375 \\
Accuracy & 0.646 & 0.563 & $+0.083$ & $[-0.010, +0.182]$ & 0.103 \\
\bottomrule
\end{tabular}
\end{center}

\noindent
No significant effect in either direction ($n=12$ paired runs). O2 is unsupported by this
run.

\subsection{E5 / E6 --- missing-modality conditions}

\begin{table}[h]\centering\small
\caption{Macro F1 by modality condition, evaluated on the 64 complete recordings.}
\begin{tabular}{llccc}
\toprule
Variant & Condition & \#\,Missing & Macro F1 & 95\,\% CI \\
\midrule
static & full          & 0 & 0.452 & [0.388, 0.518] \\
static & \code{no\_audio}    & 1 & 0.457 & [0.392, 0.527] \\
static & \code{no\_physio}   & 1 & 0.431 & [0.364, 0.495] \\
static & \code{no\_video}    & 1 & 0.459 & [0.409, 0.511] \\
static & \code{audio\_only}  & 2 & 0.444 & [0.390, 0.506] \\
static & \code{physio\_only} & 2 & 0.494 & [0.434, 0.561] \\
static & \code{video\_only}  & 2 & 0.355 & [0.259, 0.454] \\
\midrule
temporal & full          & 0 & 0.417 & [0.389, 0.450] \\
temporal & \code{no\_audio}    & 1 & 0.448 & [0.375, 0.524] \\
temporal & \code{no\_physio}   & 1 & 0.445 & [0.391, 0.512] \\
temporal & \code{no\_video}    & 1 & 0.425 & [0.389, 0.468] \\
temporal & \code{audio\_only}  & 2 & 0.443 & [0.389, 0.504] \\
temporal & \code{physio\_only} & 2 & 0.412 & [0.373, 0.451] \\
temporal & \code{video\_only}  & 2 & 0.488 & [0.428, 0.552] \\
\bottomrule
\end{tabular}
\end{table}

\noindent
The static model's best single condition is \code{physio\_only} ($0.494$), and the
temporal model's is \code{video\_only} ($0.488$) --- both \emph{above} their full-modality
scores. Removing inputs does not hurt, which is the signature of a model that is not using
them.

\subsection{E12 --- degradation curves (intended headline experiment)}

\begin{center}\small
\begin{tabular}{lcccc}
\toprule
Variant & \#\,Missing & Macro F1 & 95\,\% CI & Relative drop vs.\ full \\
\midrule
static   & 0 & 0.452 & [0.388, 0.518] & \phantom{$-$}0.0\,\% \\
static   & 1 & 0.449 & [0.413, 0.485] & \phantom{$-$}0.7\,\% \\
static   & 2 & 0.431 & [0.384, 0.479] & \phantom{$-$}4.6\,\% \\
temporal & 0 & 0.417 & [0.389, 0.450] & \phantom{$-$}0.0\,\% \\
temporal & 1 & 0.440 & [0.406, 0.476] & $-5.4$\,\% \\
temporal & 2 & 0.447 & [0.416, 0.481] & $-7.3$\,\% \\
\bottomrule
\end{tabular}
\end{center}

\noindent
The temporal model's macro F1 \emph{rises} monotonically as modalities are removed
($0.417 \rightarrow 0.440 \rightarrow 0.447$). A negative relative drop is not graceful
degradation; it indicates the removed modalities were contributing noise rather than
signal. \alert{E12 cannot be claimed from this run.}

\begin{table}[h]\centering\small
\caption{E12 paired temporal-vs-static tests per condition ($n=12$ pairs each).}
\begin{tabular}{lcccc}
\toprule
Condition & Temporal & Static & $\Delta$ & $p$ \\
\midrule
full                & 0.417 & 0.452 & $-0.035$ & 0.375 \\
\code{no\_physio}   & 0.445 & 0.431 & $+0.015$ & 0.747 \\
\code{no\_audio}    & 0.448 & 0.457 & $-0.008$ & 0.864 \\
\code{no\_video}    & 0.425 & 0.459 & $-0.034$ & 0.027 \\
\code{physio\_only} & 0.412 & 0.494 & $-0.083$ & 0.064 \\
\code{audio\_only}  & 0.443 & 0.444 & $-0.001$ & 0.973 \\
\code{video\_only}  & 0.488 & 0.355 & $+0.133$ & 0.035 \\
\bottomrule
\end{tabular}
\end{table}

\noindent
Two conditions reach nominal significance, but with seven comparisons and no multiplicity
correction, and with both models at chance level, neither should be interpreted.

\subsection{E13 --- calibration under missing modalities}

\begin{center}\small
\begin{tabular}{lccccc}
\toprule
Variant & Condition & \#\,Missing & ECE & Accuracy & Mean confidence \\
\midrule
static & full          & 0 & 0.233 & 0.563 & 0.774 \\
static & \code{no\_audio}    & 1 & 0.223 & 0.542 & 0.765 \\
static & \code{no\_physio}   & 1 & 0.263 & 0.521 & 0.782 \\
static & \code{no\_video}    & 1 & 0.150 & 0.635 & 0.768 \\
static & \code{audio\_only}  & 2 & 0.184 & 0.620 & 0.780 \\
static & \code{physio\_only} & 2 & 0.194 & 0.620 & 0.769 \\
static & \code{video\_only}  & 2 & 0.391 & 0.406 & 0.789 \\
\midrule
temporal & full          & 0 & 0.127 & 0.646 & 0.705 \\
temporal & \code{no\_audio}    & 1 & 0.192 & 0.589 & 0.697 \\
temporal & \code{no\_physio}   & 1 & 0.125 & 0.662 & 0.705 \\
temporal & \code{no\_video}    & 1 & 0.112 & 0.667 & 0.721 \\
temporal & \code{audio\_only}  & 2 & 0.097 & 0.672 & 0.721 \\
temporal & \code{physio\_only} & 2 & 0.137 & 0.620 & 0.694 \\
temporal & \code{video\_only}  & 2 & 0.142 & 0.609 & 0.685 \\
\bottomrule
\end{tabular}
\end{center}

\noindent
This is the one metric that behaves sensibly. The temporal model is markedly better
calibrated than the static model (mean ECE $0.127$ vs.\ $0.233$ at full modality) and
stays flat as modalities are removed. With accuracy at chance, however, this largely says
the temporal model is \emph{honestly} uncertain rather than usefully confident. The static
model's \code{video\_only} ECE of $0.391$ --- 79\,\% mean confidence at 41\,\% accuracy ---
is the clearest case of miscalibration in the run.

\subsection{E8 --- 3-class affect, per-class recall}

\begin{center}\small
\begin{tabular}{lccc}
\toprule
Metric & Static & Temporal & 95\,\% CI (temporal) \\
\midrule
Recall, class 0 & 0.083 & 0.083 & [0.000, 0.208] \\
Recall, class 1 & 0.240 & 0.295 & [0.146, 0.438] \\
Recall, class 2 & 0.495 & 0.449 & [0.267, 0.644] \\
Macro F1        & 0.237 & 0.233 & [0.198, 0.273] \\
\bottomrule
\end{tabular}
\end{center}

\noindent
Severe minority-class collapse --- exactly the failure mode objective O6 targets, in its
most extreme form. Class 0 recall of $0.083$ means the model almost never predicts the
minority class.

\subsection{Figures}

\IfFileExists{results/small/figures.png}{%
  \begin{figure}[h]\centering
  \includegraphics[width=\textwidth]{results/small/figures.png}
  \caption{Left: E12 degradation curves --- both flat, temporal rising. Middle: E13 ECE
  versus number of missing modalities. Right: E13 reliability diagram for the temporal
  model, full versus \code{physio\_only}.}
  \end{figure}}{%
  \begin{center}\itshape
  [Figure omitted: \code{results/small/figures.png} not found relative to this .tex file.
   Compile from the \code{research\_way/} directory, or adjust the path.]
  \end{center}}

% ===========================================================================
\section{Interpretation}

\begin{enumerate}\itemsep4pt
\item \textbf{The pipeline is correct; the result is null.} All six implemented stages run
      end to end, the splits are provably leakage-free, and the evaluation suite produces
      every intended table. The models trained successfully. They simply did not learn
      stress signal once the availability shortcut was removed.

\item \textbf{The availability confound is the dominant effect and is structural.} It is a
      property of the StressID protocol, not of this subset: audio was recorded only for
      speech tasks, which are mostly the stressful ones. Any architecture that observes
      which modalities are present --- which a missing-modality model does by construction
      --- can exploit it. Reporting a missing-modality robustness result without
      controlling for this would overstate the contribution.

\item \textbf{Scale is the most likely cause of the null result.} 96 recordings,
      $\approx$ 57 training recordings per fold, 1.2\,M parameters, encoders trained cold.
      This is exactly the small-$N$ problem Stage $-1$ was designed to solve.

\item \textbf{A larger run will not fix the confound.} It will produce better-looking
      numbers with the same defect, because the shortcut scales with the data.
\end{enumerate}

\subsection*{Recommended next steps, in order}
\begin{enumerate}\itemsep2pt
\item Decide the confound policy. The clean option is to restrict the corpus to the seven
      speech tasks, where audio is always present and availability carries zero label
      information --- at the cost of a more skewed label balance ($\approx 0.68$ positive)
      and the loss of naturally missing-modality cases.
\item Re-run at full scale (64 subjects, 11 tasks, 700 recordings) and re-measure the
      leakage inflation for E1, which is too noisy at 16 subjects to report.
\item Implement Stage $-1$ pretraining before drawing any conclusion about E2 or E12.
\item Report the \code{availability\_only} probe alongside every baseline table in the
      paper. It is cheap, and it pre-empts the most obvious reviewer objection.
\end{enumerate}

% ===========================================================================
\section{Reproducibility}

\begin{Verbatim}[fontsize=\small,frame=single,framesep=3mm]
python run_small.py                 # full run: subset -> ... -> evaluate
python run_small.py --smoke         # ~2 min sanity run
python run_small.py --stage evaluate

Environment:  Python 3.11, torch 2.6.0+cu124, CUDA on Quadro P1000 (4 GB)
              numpy 2.4.6, pandas 2.3.3, scikit-learn 1.8.0, scipy 1.17.0,
              opencv 4.13.0, matplotlib 3.10.8
Seeds:        subset/split seed 1337; training seeds 0, 1, 2
Runtime:      preprocessing ~20 min; training 1173 s; evaluation ~3 s
\end{Verbatim}

\noindent
Fixed splits are written to \code{data\_small/splits\_small.json} and reused by every
experiment. Outputs are written to \code{results/small/}: \code{baselines.csv},
\code{e1\_leakage\_inflation.csv}, \code{shortcut\_audit.csv}, \code{predictions.csv},
\code{per\_run\_metrics.csv}, \code{e2\_temporal\_vs\_static.csv},
\code{e5e6\_missing\_modality.csv}, \code{e8\_per\_class\_recall.csv},
\code{e12\_degradation.csv}, \code{e12\_paired\_tests.csv}, \code{e13\_calibration.csv},
\code{e13\_reliability\_bins.csv}, \code{report.json}, \code{figures.png}.

% ===========================================================================
\section{Deviations from the planning document}

\begin{enumerate}\itemsep3pt
\item \textbf{10\,s windows, not 30\,s.} The plan assumes 3--5 minute recordings. Real
      StressID speech tasks are 60\,s, so 30\,s windows at 50\,\% overlap would yield only
      three windows per task --- too short a sequence to test temporal modelling at all.
      10\,s windows with a 5\,s hop give 11 windows.

\item \textbf{Log-mel + CNN instead of wav2vec2.} The \code{transformers} package is not
      installed and a wav2vec2 fine-tune does not fit a 4\,GB GPU at this scale. The audio
      encoder is an isolated drop-in slot; substituting wav2vec2 makes E9 runnable without
      touching the rest of the model.

\item \textbf{Face-crop pixels instead of OpenFace action units.} OpenFace is unavailable.
      Video is decoded at 2.5\,fps, the face is located once per clip with a Haar cascade,
      and the crop is resized to $32\times32$ grayscale. The video encoder is likewise a
      drop-in slot for per-frame AU vectors. The repository ships \emph{task-level} AU
      aggregates, but those are static per-task vectors --- precisely the representation
      the temporal stage exists to avoid.

\item \textbf{Static baseline shares all weights-bearing components with the temporal
      model.} The two variants differ only in Stage 4 (mean pooling versus a Transformer
      over time), so the E12 ablation isolates temporal modelling and nothing else.

\item \textbf{Window cap.} \code{max\_windows = 16} truncates Breathing (180\,s) and Relax
      (150\,s) to their first $\approx 85$\,s. E5/E6/E12 are unaffected, being scored on
      the 60\,s speech tasks only.
\end{enumerate}

\vspace{1em}
\noindent\rule{\textwidth}{0.4pt}\\[2pt]
\small\textit{Generated from the run in \code{research\_way/results/small/}. All numbers
in this report are read directly from those files; none were transcribed by hand from
intermediate output.}

\end{document}
